\documentclass[a4paper,12pt]{article}

% Pacchetti fondamentali
\usepackage[utf8]{inputenc}   % Codifica caratteri
\usepackage[italian]{babel}    % Lingua italiana
\usepackage{amsmath}           % Formule matematiche avanzate
\usepackage{graphicx}          % Inserimento immagini
\usepackage{geometry}          % Gestione margini
 \geometry{margin=2.5cm}
\usepackage{booktabs}          % Tabelle professionali
\usepackage{physics}
\usepackage{float}
\usepackage{caption}
\usepackage{textcomp}
\usepackage{amssymb}
\usepackage{parskip}

\title{Misura della costante elastica di una molla elicoidale}
\author{Eddi Bressan}
\date{\today}

\begin{document}

\maketitle
\tableofcontents
\newpage

\section{Obiettivo dell'esperienza}
Lo scopo dell'esperienza è lo studio della deformazione di una molla elicoidale tramite 2 metodi: statico e dinamico.

\section{Cenni teorici}

\subsection{Deformazione elastica}
Una deformazione è detta \textbf{elastica} se il corpo torna allo stato originario quando
vengono meno le forze che ne hanno causato la deformazione.

Tale deformazione può essere definita \textbf{elastica} se le forze applicate sono inferiori
ad un limite dipendente da vari fattori: materiale, temperatura, tipo di deformazione considerata, etc.
Oltre questo limite le deformazioni divengono anaelastiche, ovvero permanenti.

Nelle deformazioni elastiche si osserva una relazione di proporzionalità tra sollecitazione e deformazione:
$$
\Delta L \propto F
$$
\noindent
Questo comportamento è noto come \textbf{legge di Hooke}.

La legge di Hooke è valida per la maggior parte dei minerali, per il vetro,
per i materiali ceramici e per i metalli. Per i materiali duttili invece è vera per carichi modesti.

Quando all'estremità libera di una molla è applicata una forza,
la molla esercita una forza di richiamo proporzionale allo spostamento di tale 
estremo rispetto alla posizione di equilibrio.

Il modulo della forza di richiamo è proporzionale alla deformazione (legge di Hooke):
$$
\vec{F} = -k \vec{x}
$$
Dove $\vec{x}$ è la deformazione della molla rispetto la sua lunghezza normale $\vec{x} = (l - l_{0})\hat{x}$
con $l_0$ la lunghezza della molla a riposo.

La costante di proporzionalità $k$ è detta \textbf{costante elastica} della molla e 
dipende dal materiale di cui è costituita, dal diametro del filo e dal numero di spire della molla.


\subsection{Metodo statico}
La costante elastica di una molla può essere determinata sperimentalmente misurando
le elongazioni $\Delta l$ al variare delle forze applicate.

Se all'estremità inferiore di una molla sospesa verticalmente è appesa una massa m,
la forza applicata coincide con la forza peso della massa e la legge di Hooke si riscrive
nel seguente modo:
$$
mg = k {\Delta l} = k (l-l_0)
$$
Usando diverse masse e misurando i rispettivi $\Delta l$ ottengo diverse misure di k:
$$
k_i = \frac{m_i g}{\Delta l_i}
$$
Nel caso reale alla molla è già appeso un piattello (di massa $m_p$),
inoltre la molla ha una sua massa ed è quindi anch'essa soggetta alla forza peso.
Subisce già un allungamento $l_s$ anche a piattello scarico.

Misurando la differenza tra l'allungamento a piatto scarico:
$$
(m_p + m_{eff.statica})g = k(l_s - l_0)
$$
e l'allungamento in seguito all'aggiunta di una massa $m_i$:
$$
(m_i + m_p + m_{eff.statica})g = k(l_i - l_0)
$$
e facendo la differenza membro a membro tra le due operazioni otteniamo:
$$
m_i g = k_i (l_i - l_s)
$$
Da cui il valore della costante elastica della molla:
$$
k_i = \frac{m_i g}{(l_i - l_s)}
$$
Gli errori sono tutti di sensibilità per cui possiamo scrivere:
$$
\frac{\Delta k_i}{k_i} = \frac{\Delta m_i}{m_i}+2\frac{\Delta l}{(l_i - l_s)}
$$

\subsection{Metodo dinamico}
Partendo dalla legge di Hooke:
$$
mg = k \Delta l = k(l_{eq} - l_0)
$$
Se applicando una forza alla massa appesa alla molla sposto la molla dalla sua posizione di equilibrio e la lascio
andare, questa si metterà ad oscillare rispetto la posizione di equilibrio $l_{eq}$ con equazione:
$$
m\ddot{x} + kx = mg
$$
La soluzione di questa equazione differenziale sarà data da una soluzione dell'equazione
omogenea:
$$
m\ddot{x} + kx = 0 \quad \text{o} \quad \ddot{x} + \frac{k}{m}x = 0
$$
Sommata ad una soluzione particolare dell'equazione non omogenea.
Le condizioni iniziali sono:
$$
\begin{cases}
    \dot{x}(t=0)=0 \\
    x(t=0)=x_0
\end{cases}
$$
La soluzione dell'equazione differenziale omogenea è del tipo $x(t) = A \sin(\omega t+\phi)$
mentre una soluzione particolare dell'equazione non omogenea è la soluzione costante:
$$
x = \frac{mg}{k} = l_{eq} - l_0
$$
La soluzione completa è quindi:
$$
x(t) = A \sin(\omega t+\phi) + l_{eq} - l_0
$$
Con:
$$
\begin{cases}
    \dot{x}(t=0)=0 \\
    x(t=0)=x_0
\end{cases}
$$
Abbiamo quindi:
$$
\begin{cases}
    \dot{x} = A \omega \cos(\omega t + \phi) = 0 \\
    \ddot{x} = -A \omega^2 \sin(\omega t + \phi)
\end{cases}
$$
Da cui ricaviamo:
$$
-A \omega^2 \sin(\omega t + \phi) + \frac{k}{m} A \sin(\omega t + \phi) + \frac{k}{m}[l_{eq} - l_0] = g
$$
Dobbiamo perciò avere:
$$
\omega = \sqrt{\frac{k}{m}}
$$
Ovvero la pulsazione non cambia a causa della forza di gravità.

La condizione iniziale $\dot{x}(t=0)=0$ fornisce:
$$
\cos(\phi) = 0 \quad \text{o} \quad \phi = \frac{\pi}{2} 
$$
Mentre applicando la condizione $x(t=0)=x_0$ otteniamo:
$$
x_0 = A + l_{eq} - l_0
$$
Da cui:
$$
A = x_0 - (l_{eq} - l_0)
$$
Per una soluzione completa data da:
$$
x(t) = [x_0 - (l_{eq} - l_0)]\sin(\sqrt{\frac{k}{m}}t+\frac{\pi}{2}) + (l_{eq} - l_0) 
$$
Il fatto che $\omega = \sqrt{\frac{k}{m}}$ ci permette di calcolare la costante elastica della molla dalla
misura del periodo di oscillazione:
$$
T = \frac{2\pi}{\omega} = 2\pi \sqrt{\frac{m}{k}}
$$
Ma in questo caso la massa è tutta la massa appesa alla molla, inclusa la massa del supporto.

\subsection{Massa efficace statica}
Per calcolare la $m{eff.statica}$ immaginiamo di dividere la molla in N parti uguale lunghezza.
La lunghezza di tali sezioni sarà $a = l_0 / N$ e la massa pari a $\mu = m_m / N$.
La costante elastica della molla completa è k, quindi, considerando tale molla come una serie di N molle,
la costante elastica k di ognuna di queste parti della molla sarà data da:
$$
k = \frac{1}{\sum\frac{1}{k}} = \frac{k}{N}
$$
Il primo elemento della molla sentirà il peso di tutte le N-1 molle e si allungherà di:
$$
\Delta x_1 = \frac{(N-1)\mu g}{k} 
$$
La seconda sentirà il peso delle N-2 molle e si allungherà di:
$$
\Delta x_2 = \frac{(N-2)\mu g}{k} 
$$
L'allungamento totale sarà dato dalla somma di tutti gli allungamenti:
$$
\Delta x = \sum_{i=1}^{N-1} \Delta x_i = \sum_{i=1}^{N-1} \frac{(N-i)\mu g}{k} \implies \Delta x = \frac{m_m g}{2} \frac{1}{k}
$$
Ovvero la molla si è allungata come se fosse stata appesa una massa:
$$
m_{eff.statica} = \frac{m_m}{2}
$$

\subsection{Massa efficace dinamica}
Nella misura tramite metodo dinamico non possiamo trascurare la massa della molla in relazione alla stima del
periodo di oscillazione.
Dovremo valutare il contributo di $m_m$ anche in questo caso come $m_{eff.dinamica}$

Calcoliamo la variazione dell'energia cinetica $dK$ della molla a causa della stessa massa $dm$ di un 
elemento di lunghezza $dx$ che si muove con una velocità istantanea $v$ durante il moto armonico della molla:
$$
dK = \frac{1}{2}dmv^2
$$
Poiché vale $dm = m_m \frac{dx}{l}$, mentre la velocità alla distanza x sarà legata alla velocità all'estremo
$l$ della molla V da:
$$
v = \frac{x}{l}V
$$
Possiamo quindi scrivere:
$$
dK = \frac{1}{2}(\frac{x}{l}V)^2 \frac{m_m}{l}dx
$$
L'energia cinetica totale della molla in oscillazione si ottiene integrando da 0 alla lunghezza della molla $l$:
$$
K = \frac{1}{2}(\frac{V}{l})^2 \frac{m_m}{l} \int_{0}^{l} x^2 dx
$$
Che vale:
$$
K = \frac{1}{2}(\frac{V}{l})^2 \frac{m_m}{l} \frac{l^3}{3} = \frac{1}{2}V^2 \frac{m_m}{3}
$$
Sommando anche la massa appesa alla molla (anch'essa in movimento con velocità $V$) abbiamo un'energia
cinetica totale:
$$
K = \frac{1}{2}V^2 (m+\frac{m_m}{3})
$$
Abbiamo quindi che la massa effettiva della molla da sommare alla massa appesa è nel caso dinamico pari a:
$$
m_{eff.dinamica} = \frac{m_m}{3}
$$

\subsection{Metodo dinamico con attriti}
Nel caso in cui non si possano trascurare gli attriti, dobbiamo inserire un termine proporzionale al
$dx/dt$ che dipenderà daòòa viscosità del mezzo (e da eventuali attriti meccanici nella molla stessa).
Consideriamo per semplicità l'equazione omogenea (abbiamo visto la soluzione della non omogenea):
$$
m\ddot{x} + C\dot{x} + kx = 0
$$
Ovvero:
$$
\ddot{x} + \frac{C}{m} \dot{x} + \frac{k}{m} x = 0
$$
Con condizioni iniziali:
$$
x(t=0) = x_0 \quad \text{e} \quad \dot{x}(t=0) = 0
$$
E con, per comodità:
$$
2\gamma = \frac{C}{m} \quad \text{e} \quad \omega_0 = \sqrt{\frac{k}{m}}
$$
Ovvero:
$$
\ddot{x} + 2\gamma \dot{x} + \omega_0^2 x = 0
$$
Che è un'equazione differenziale lineare di secondo ordine a coefficienti costanti.

Dalla teoria delle equazioni differenziali ricavo il moto descritto dalla legge:
$$
x(t) = A e^{\gamma t} \cos(\sqrt{\omega_0^2 - \gamma^2 t} + \phi) \quad \text{per} \quad \gamma < \omega_0
$$
Ottengo quindi un moto oscillatorio smorzato, con A e $\phi$ che dipendono dalle condizioni iniziali.
Notiamo che la pulsazione misurata è influenzata dall'attrito:
$$
\omega_0 \to \omega_0 - \frac{C}{2m}
$$

\newpage
\section{Apparato sperimentale}
\textbf{Strumenti di misura}
\begin{table}[htbp]
    \centering
    \begin{tabular}{lr} % l = left (allineato a sinistra), c = center (centrato)
        \toprule
        \textbf{Strumento} & \textbf{Risoluzione} \\
        \midrule
        Metro a nastro      & $1$ mm               \\
        Cronometro digitale & $0.01$ s             \\
        Bilancia analitica    & $0.01$ g              \\
        \bottomrule
    \end{tabular}
\end{table}\\
\textbf{Componenti del sistema fisico}
\begin{table}[htbp]
    \centering
    \footnotesize
    \vspace{8pt} % Spazio aggiuntivo tra didascalia e tabella
    \begin{tabular}{lll}
        \toprule
        \textbf{Elemento} & \textbf{Proprietà} & \textbf{Valore Nominale} \\
        \midrule
        Massa oscillante (Piattello) & Massa ($M_p$)     & 20.10 g                \\
        Massa appesa                 & Massa ($M_i$)  & variabile (4.97 - 49.99)g \\
        \addlinespace
        Molla      & Massa ($M_m$)     &  29.96 g                  \\
                                 & Lunghezza ($l_M$) & variabile (19.9 - 65.3) cm           \\
        \addlinespace
        Stativo                  & Supporto rigido   & ---                    \\
        \bottomrule
    \end{tabular}
\end{table}

\section{Procedura sperimentale}
\subsection{Misure preliminari}
\begin{itemize}
    \item Misura della massa della molla $M_m$, del piattello $M_p$ e delle masse $M_i$.
    \item Misura della lunghezza della molla a riposo $l_0$.
\end{itemize}

\subsection{Set-up dell'apparato}

\end{document}